\documentclass[12pt]{article}
\usepackage{amsmath, amssymb}
\usepackage[margin=1in]{geometry}
\setlength{\parskip}{6pt}
\title{QUBO Formulation: Quantum Circuit Compilation Candidate Selection Problem}
\date{}
\begin{document}
\maketitle

\subsection*{Quantum Circuit Compilation Candidate Selection Problem}

\paragraph{Problem Statement.}
Given a set of $N$ candidate compiled circuits, the objective is to select exactly one candidate that minimizes the total compilation cost. Each candidate $i \in \{0, \dots, N-1\}$ is associated with a pre-computed scalar cost $c_i$, which aggregates depth, fidelity loss, and hardware constraint penalties. The problem requires choosing a single candidate such that the sum of the costs of the selected candidates is minimized, subject to the constraint that precisely one candidate is chosen.

\paragraph{Decision Variables.}
We introduce binary decision variables $x_i \in \{0, 1\}$ for each candidate $i \in \{0, \dots, N-1\}$. The variable $x_i$ takes the value $1$ if candidate $i$ is selected, and $0$ otherwise.

\paragraph{QUBO Derivation.}
We formulate the problem as a Quadratic Unconstrained Binary Optimization (QUBO) by converting the constrained minimization into an unconstrained energy minimization.

\textbf{Step 1 --- Objective Function.}
The original objective is to minimize the total cost:
\begin{equation}
    \min_{x} \; C(x) = \sum_{i=0}^{N-1} c_i x_i.
\end{equation}
Since QUBO minimization is equivalent to the original minimization, we retain this form.

\textbf{Step 2 --- Constraint.}
The requirement to select exactly one candidate is expressed as:
\begin{equation}
    \sum_{i=0}^{N-1} x_i = 1.
\end{equation}

\textbf{Step 3 --- Penalty Term.}
We enforce the constraint by adding a quadratic penalty term $P(x)$ with weight $\lambda > 0$:
\begin{equation}
    P(x) = \lambda \left( \sum_{i=0}^{N-1} x_i - 1 \right)^2.
\end{equation}
Expanding the square yields:
\begin{equation}
    P(x) = \lambda \left( \sum_{i=0}^{N-1} x_i^2 + \sum_{i \neq j} x_i x_j - 2 \sum_{i=0}^{N-1} x_i + 1 \right).
\end{equation}
Using the idempotent property of binary variables, $x_i^2 = x_i$, we simplify:
\begin{align}
    P(x) &= \lambda \left( \sum_{i=0}^{N-1} x_i + \sum_{i \neq j} x_i x_j - 2 \sum_{i=0}^{N-1} x_i + 1 \right) \\
    &= \lambda \left( -\sum_{i=0}^{N-1} x_i + \sum_{i \neq j} x_i x_j + 1 \right) \\
    &= -\lambda \sum_{i=0}^{N-1} x_i + \lambda \sum_{i \neq j} x_i x_j + \lambda.
\end{align}

\textbf{Step 4 --- Combined Hamiltonian.}
The total QUBO Hamiltonian $H(x)$ is the sum of the objective and penalty terms:
\begin{align}
    H(x) &= C(x) + P(x) \\
    &= \sum_{i=0}^{N-1} c_i x_i - \lambda \sum_{i=0}^{N-1} x_i + \lambda \sum_{i \neq j} x_i x_j + \lambda \\
    &= \sum_{i=0}^{N-1} (c_i - \lambda) x_i + \lambda \sum_{i \neq j} x_i x_j + \lambda.
\end{align}

\textbf{Step 5 --- Q Matrix and Offset.}
The Hamiltonian is expressed in the standard QUBO form $H(x) = \sum_{i \le j} Q_{ij} x_i x_j + c$. Comparing terms, the matrix entries $Q_{ij}$ and offset $c$ are given by:
\begin{align}
    Q_{ii} &= c_i - \lambda, \quad \forall i \in \{0, \dots, N-1\}, \\
    Q_{ij} &= \lambda, \quad \forall i, j \in \{0, \dots, N-1\}, \; i \neq j, \\
    c &= \lambda.
\end{align}
Note that the term $\lambda \sum_{i \neq j} x_i x_j$ corresponds to symmetric matrix entries $Q_{ij} = \lambda$ for $i \neq j$.

\paragraph{Penalty Weight Justification.}
The penalty weight $\lambda$ must be chosen sufficiently large to ensure that any infeasible solution has strictly higher energy than any feasible solution. A sufficient condition is:
\begin{equation}
    \lambda > \max_{i \in \{0, \dots, N-1\}} c_i.
\end{equation}
This choice guarantees that the minimum penalty for a constraint violation (which occurs when zero candidates are selected, yielding penalty $\lambda$) exceeds the cost of any single candidate. For violations involving multiple selections ($k \ge 2$), the penalty grows quadratically as $\lambda(k-1)^2 \ge \lambda$, further increasing the energy. Thus, $\lambda > \max_i c_i$ ensures that the global minimum of $H(x)$ cannot occur at an infeasible point.

\paragraph{Correctness Argument.}
Minimizing $H(x)$ over $\{0, 1\}^N$ recovers the optimal compiled circuit candidate. For any feasible solution $x$ (where $\sum_i x_i = 1$), the penalty term vanishes, and $H(x) = C(x)$. For any infeasible solution, $H(x) = C(x) + P(x)$. Under the condition $\lambda > \max_i c_i$, the energy of any infeasible solution exceeds $\max_i c_i$, while the energy of any feasible solution is at most $\max_i c_i$. Consequently, the global minimum must be feasible. Among feasible solutions, $H(x)$ equals the objective $C(x)$, so the minimizer of $H(x)$ is the candidate with the minimum compilation cost.

\paragraph{Illustrative Example.}
Consider an instance with $N=3$ candidates and costs $c_0 = 10$, $c_1 = 20$, $c_2 = 15$. We choose $\lambda = 21$, satisfying $\lambda > \max\{10, 20, 15\}$.

The concrete QUBO parameters are:
\begin{align}
    Q_{00} &= 10 - 21 = -11, \\
    Q_{11} &= 20 - 21 = -1, \\
    Q_{22} &= 15 - 21 = -6, \\
    Q_{ij} &= 21 \quad \text{for } i \neq j, \\
    c &= 21.
\end{align}
The Hamiltonian is:
\begin{equation}
    H(x) = -11 x_0 - x_1 - 6 x_2 + 21(x_0 x_1 + x_0 x_2 + x_1 x_2) + 21.
\end{equation}
Evaluating $H(x)$ for candidate solutions:
\begin{itemize}
    \item $x = [1, 0, 0]$: $H = -11 + 21 = 10$.
    \item $x = [0, 1, 0]$: $H = -1 + 21 = 20$.
    \item $x = [0, 0, 1]$: $H = -6 + 21 = 15$.
    \item $x = [0, 0, 0]$: $H = 21$.
    \item $x = [1, 1, 0]$: $H = -11 - 1 + 21 + 21 = 30$.
\end{itemize}
The global minimum is $H = 10$ at $x = [1, 0, 0]$, corresponding to selecting candidate $0$ with minimum cost $10$.
\end{document}